% Part: incompleteness
% Chapter: representability-in-q
% Section: prim-rec

\documentclass[../../../include/open-logic-section]{subfiles}

\begin{document}

\olfileid{inc}{req}{pri}
\olsection{模拟原始递归}

现在我们可以证明，利用 beta 函数，原始递归定义可以被正则极小化“模拟”。
假设我们有 $f(\vec x)$ 和 $g(\vec x, y, z)$。
那么由 $f$ 和~$g$ 通过原始递归定义的函数~$h(x,\vec z)$ 为
\begin{align*}
h(\vec x, 0) & =  f(\vec x) \\
h(\vec x, y+1) & =  g(\vec x, y, h(\vec x, y)).
\end{align*}
我们需要证明，$h$ 可以仅用复合和正则极小化从 $f$ 和~$g$ 定义出来，
所用的函数包括基本函数以及从它们出发通过复合和正则极小化定义的函数（例如~$\beta$）。

\begin{lem}
\ollabel{lem:prim-rec}
如果 $h$ 可以由 $f$ 和 $g$ 利用原始递归定义，
那么它也可以从 $f$、$g$ 以及函数 $\Zero$、$\Succ$、$\Proj{n}{i}$、$\Add$、$\Mult$、$\Char{=}$ 出发，
利用复合和正则极小化定义出来。
\end{lem}

\begin{proof}
首先，定义一个辅助函数 $\hat h(\vec x, y)$，它返回最小的数~$d$，使得 $d$ 编码一个满足下列条件的序列：
\begin{enumerate}
\item $(d)_0 = f(\vec x)$，并且
\item 对每个 $i < y$，$(d)_{i+1} = g(\vec x, i, (d)_i)$，
\end{enumerate}
其中 $(d)_i$ 是 $\beta(d,i)$ 的简写。换言之，$\hat h$ 返回序列 $\tuple{h(\vec x, 0), h(\vec x, 1), \dots, h(\vec
x, y)}$。
我们可以将 $\hat h$ 写成
\[
\hat h(\vec x, y) = \umin{d}{(\beta(d,0) = f(\vec x) \land \bforall{i <
  y}{\beta(d,i+1) = g(\vec x, i,\beta(d,i)})}.
\]
注意：这里不需要原始递归，只需要极小化。
根据 beta 函数引理 \olref[bet]{lem:beta}，被极小化的函数是正则的。

但现在我们有
\[
h(\vec x, y) = \beta(\hat h(\vec x, y), y),
\]
因此，$h$ 可以仅用复合和正则极小化从基本函数定义出来。
\end{proof}

\end{document}